\small
\textbf{Abstract.}
Which forecasting model is best depends on the level of the product and store
hierarchy at which the forecast is judged, and on how much of the sales the
plan (price, promotion, events) moves. We compare a world model, which carries
each item's state forward one period at a time under the plan, with a direct
lift model, LightGBM, a graph model, seasonal naive, and the pretrained
TimesFM~2.5 with and without the plan, on three public data sets: M5 (four
28-day windows), six Dominick's categories and Breakfast at the Frat. Every
model is scored at every level of the hierarchy, and every model and change
that lost is reported. On M5, TimesFM given only each item's sales history is
the most accurate model item by store in all four months
(\mTfmMarIxs{} to \mTfmTestIxs{}), but forecasts \tfmUnderLo\% to
\tfmUnderHi\% below total sales; the world model is the most accurate at the
chain total in two of the four months and LightGBM in one. The official M5 test
month is the only one of the four in which LightGBM is the most accurate over
all levels, so a single test month would have given the wrong ranking. On the
weekly, promotion-driven data, TimesFM without the plan forecasts
\promoTfmLo\% to \promoTfmHi\% of promoted-week sales and is behind both
plan-aware trained models item by store in six of seven data sets; given the
plan, it forecasts \promoTfmPlanLo\% to \promoTfmPlanHi\% and becomes the most
accurate model in one category. For an item whose first price cut in a year
sold \cutSold{} units, the world model forecast \cutWorld{} across three seeds,
against about 3 for LightGBM and 8 for TimesFM: plan-aware models respond in
the right direction, and the data do not fix the size. A retailer should
choose its forecasting model by the decision the forecast feeds, scored at
that decision's level and over several periods.

\vspace{4pt}
\textbf{Keywords:} retail forecasting; hierarchical evaluation; promotions;
world models; foundation models; M5.
